\documentclass[11pt,a4paper]{article}
\usepackage[utf8]{inputenc}
\usepackage[T1]{fontenc}
\usepackage{lmodern}
\usepackage[french]{babel}
\usepackage[margin=1.9cm]{geometry}
\usepackage{microtype}
\usepackage{booktabs}
\usepackage{array}
\usepackage[table]{xcolor}
\usepackage{enumitem}
\usepackage[most]{tcolorbox}
\definecolor{navy}{HTML}{1F3864}
\definecolor{bluemid}{HTML}{2E5496}
\definecolor{lightblue}{HTML}{D6E0F0}
\definecolor{accent}{HTML}{C0491F}
\newtcolorbox{cle}[1][]{colback=lightblue!40,colframe=navy,boxrule=0.6pt,
  arc=2pt,left=7pt,right=7pt,top=5pt,bottom=5pt,before skip=6pt,after skip=6pt,#1}
\setlength{\parindent}{0pt}
\setlength{\parskip}{4pt}

\title{\vspace{-1.2cm}\color{navy}\bfseries Dix incidents, une question simple\\[2pt]
  \large Dans chaque incident, qu'est-ce qui a lâché en premier ?}
\author{Kélian Kaddouri}
\date{}

\begin{document}
\maketitle
\vspace{-1.1cm}

\begin{cle}
\textbf{Ce qu'on vous demande.} Vous allez lire dix récits d'incidents réels, tirés de rapports
d'enquête publics. Pour chacun, on vous demande simplement, entre deux domaines qui ont tous les
deux flanché, \textbf{lequel a flanché en premier et a entraîné l'autre}. C'est tout. Comptez
une heure. Aucune connaissance technique n'est nécessaire : les faits vous sont
donnés, vous n'avez qu'à les ordonner.
\end{cle}

\section*{Les cinq domaines}

Toute la grille porte sur ces cinq domaines. Retenez les noms, ils suffisent.

\begin{center}
\renewcommand{\arraystretch}{1.35}
\begin{tabular}{@{}>{\bfseries}p{3.1cm} p{13.2cm}@{}}
\toprule
Gouvernance & Qui décide, qui est responsable de quoi, quelle importance la direction accorde à
la sécurité, quels moyens elle y met. \\
Incidents & Repérer qu'un incident a lieu, l'escalader, y répondre, en informer les clients et
les autorités. \\
Tests & Tester avant de mettre en service, vérifier régulièrement, appliquer les correctifs,
renouveler les certificats et les clés. \\
Prestataires & Les fournisseurs et sous-traitants, et la dépendance à un outil ou un service
extérieur. \\
Partage d'infos & Recevoir et faire circuler les alertes sur les menaces, en interne comme avec
l'extérieur. \\
\bottomrule
\end{tabular}
\end{center}

\section*{Comment répondre}

Pour chaque case, vous avez le nom de deux domaines. Écrivez l'une de ces quatre réponses :

\begin{center}
\renewcommand{\arraystretch}{1.3}
\begin{tabular}{@{}>{\bfseries}p{4.6cm} p{11.6cm}@{}}
\toprule
le nom d'un domaine & si le récit montre que celui-là a lâché en premier \emph{et} a entraîné
l'autre. Écrivez simplement son nom, par exemple \emph{Gouvernance}. \\
en même temps & les deux ont lâché, mais rien ne dit lequel a entraîné l'autre. \\
non concerné & au moins un des deux domaines n'apparaît pas dans ce récit. \\
je ne sais pas & vous hésitez, ou le récit ne permet pas de trancher. \\
\bottomrule
\end{tabular}
\end{center}

\begin{cle}[colback=accent!7,colframe=accent!70!black]
\textbf{La règle à tenir, et la seule.} N'écrivez le nom d'un domaine que si le récit
\textbf{montre} qu'il a entraîné l'autre. Si les deux problèmes sont juste racontés l'un à côté
de l'autre, répondez \emph{en même temps}. Dans le doute, \emph{je ne sais pas} est une bonne
réponse. Il n'y a pas de réponse attendue, et si vous répondez souvent \emph{en même temps} ce
n'est pas un problème : c'est même une information.
\end{cle}

\newpage
\section*{Les sept récits}

\subsection*{1. Fournisseur de messagerie d'entreprise, 2023}
\textit{Rapport d'une commission d'enquête publique américaine.}
\begin{itemize}[itemsep=2pt]
\item Le rapport juge que l'entreprise ne mettait pas assez la sécurité au premier plan, et
      recommande que la direction générale et le conseil s'en saisissent eux-mêmes.
\item Le renouvellement des clés de sécurité, fait à la main et rarement, a été arrêté en 2021
      après une panne, sans être remplacé par un système automatique.
\item Aucune alerte ne signalait qu'une clé devenait trop ancienne.
\item Une clé créée en 2016, censée ne plus servir, a été utilisée par les attaquants.
\item L'entreprise n'a pas découvert l'attaque elle-même : c'est un client qui l'a prévenue.
\item Ses communications publiques donnaient une explication de l'origine de la fuite ; elles ont
      été corrigées six mois plus tard pour dire que l'origine restait inconnue.
\item Le rapport relève des failles dans la façon dont l'entreprise sécurise les sociétés qu'elle
      rachète : les accès d'un ingénieur venu d'une société rachetée ont servi à entrer.
\end{itemize}

\subsection*{2. Banque britannique, migration informatique de 2018}
\textit{Enquête indépendante commandée par la banque ; sanctions des deux autorités de contrôle
britanniques ; sanction personnelle de l'ancien directeur informatique.}
\begin{itemize}[itemsep=2pt]
\item Cinq millions de clients ont été basculés sur une nouvelle plateforme. Beaucoup se sont
      retrouvés sans accès à leur compte, certains ont vu des montants disparaître, d'autres ont
      eu accès aux comptes d'autres clients.
\item Les autorités ont sanctionné des manquements dans le pilotage du projet et dans la
      surveillance des sous-traitants.
\item Le chantier a mobilisé plus de 70 fournisseurs et plus de 1\,400 personnes ; l'enquête juge
      la surveillance de ces fournisseurs insuffisante.
\item Deux centres de données n'ont pas été testés avant la mise en service, l'un étant réservé à
      l'exploitation : l'ensemble n'a donc jamais été testé en charge réelle.
\item Sur 5\,359 défauts recensés, 4\,424 étaient encore ouverts le jour du démarrage.
\item L'ancien directeur informatique a été sanctionné à titre personnel pour n'avoir pas veillé
      à ce que la banque encadre convenablement son contrat de sous-traitance.
\item La gestion de la crise et l'information des clients ont été critiquées ; le coût dépasse
      300 millions de livres.
\end{itemize}

\subsection*{3. Société américaine de renseignement financier, 2017}
\textit{Rapports d'un organisme d'audit public et d'une commission parlementaire.}
\begin{itemize}[itemsep=2pt]
\item En mars 2017, une alerte publique demande de corriger sans délai un logiciel très répandu.
      L'entreprise diffuse l'alerte en interne, mais la liste de diffusion n'était plus à jour et
      les personnes concernées ne l'ont jamais reçue.
\item Un contrôle automatique passé la semaine suivante n'a pas repéré la faille.
\item Un certificat de surveillance avait expiré en janvier 2017 : sans lui, les outils de
      surveillance ne voyaient plus rien. Personne ne s'en est aperçu pendant plus de six mois.
      Plus de 300 certificats étaient expirés, dont 79 sur des systèmes critiques.
\item Le certificat a été renouvelé fin juillet 2017 : les outils ont immédiatement signalé une
      activité suspecte. Les intrus étaient là depuis environ 76 jours.
\item L'audit relève qu'il n'y avait pas de responsabilités claires ni de ligne hiérarchique nette
      dans l'organisation informatique, et un écart entre les règles écrites et leur application
      réelle, qui a aussi freiné d'autres chantiers de sécurité.
\item L'annonce publique a eu lieu en septembre 2017.
\end{itemize}

\subsection*{4. Société de courtage américaine, 1\up{er} août 2012}
\textit{Décision du régulateur boursier américain.}
\begin{itemize}[itemsep=2pt]
\item Une mise à jour logicielle n'a pas été installée correctement sur tous les serveurs ; l'un
      d'eux a continué à faire tourner un ancien programme resté actif.
\item Le régulateur constate l'absence de contrôles et de procédures convenables pour installer
      et vérifier ces mises à jour.
\item Il constate aussi l'absence de procédures écrites suffisantes pour guider les employés
      quand un incident technique grave survient.
\item L'épisode a duré environ 45 minutes et coûté de l'ordre de 440 millions de dollars.
\end{itemize}

\subsection*{5. Banque américaine, 2019}
\textit{Décisions de deux autorités de contrôle bancaires américaines.}
\begin{itemize}[itemsep=2pt]
\item L'autorité sanctionne l'absence d'une véritable évaluation des risques \emph{avant} de
      basculer des activités importantes chez un hébergeur extérieur.
\item Elle relève aussi que les insuffisances constatées n'ont pas été corrigées dans un délai
      raisonnable.
\item Un équipement de protection mal configuré a permis d'accéder aux données hébergées.
\item La banque a appris l'intrusion par un signalement venu de l'extérieur.
\item La seconde décision porte sur le pilotage et la surveillance par le conseil
      d'administration.
\end{itemize}

\subsection*{6. Éditeur de logiciels pour les marchés financiers, janvier 2023}
\textit{Communications de l'entreprise et réactions des autorités européennes et britanniques.}
\begin{itemize}[itemsep=2pt]
\item Un rançongiciel a frappé une division de cet éditeur, dont les logiciels sont utilisés par
      de nombreux acteurs des marchés.
\item Beaucoup de ses clients ont dû revenir à des traitements manuels.
\item Des déclarations réglementaires ont été retardées dans tout le secteur.
\end{itemize}

\subsection*{7. Logiciel de transfert de fichiers, mai-juin 2023}
\textit{Listes de victimes publiques.}
\begin{itemize}[itemsep=2pt]
\item Une faille d'un logiciel de transfert de fichiers vendu par un éditeur extérieur a été
      exploitée massivement.
\item Des centaines d'organisations, dont des banques et des assureurs, ont été touchées en même
      temps et ont dû gérer puis déclarer un incident.
\end{itemize}

\subsection*{8. Éditeur d'un logiciel de sécurité pour postes de travail, juillet 2024}
\textit{Analyse technique publiée par l'entreprise elle-même.}
\begin{itemize}[itemsep=2pt]
\item Une mise à jour de contenu diffusée automatiquement contenait un fichier dont le nombre de
      champs ne correspondait pas au modèle attendu : vingt-et-un champs prévus, vingt fournis.
\item Le contrôle automatique chargé de valider ces fichiers avant diffusion ne l'a pas arrêté.
\item Les couches de test ne l'ont pas vu non plus, à cause d'un mécanisme de correspondance
      générique introduit plusieurs mois plus tôt.
\item Environ huit millions et demi de machines se sont arrêtées, dans des hôpitaux, des banques,
      des aéroports et des administrations.
\item L'entreprise a dû gérer une crise mondiale et publier une procédure de réparation manuelle,
      machine par machine.
\end{itemize}

\subsection*{9. Bibliothèque logicielle libre très répandue, fin 2021}
\textit{Rapport d'une commission d'enquête publique américaine.}
\begin{itemize}[itemsep=2pt]
\item Une faille critique est révélée dans un composant logiciel gratuit, utilisé sans le savoir
      par une très grande partie des organisations.
\item Beaucoup d'entre elles n'ont pas pu dire où ce composant était installé chez elles : elles
      ne tenaient pas d'inventaire de ce qu'elles utilisaient et n'avaient pas de règles sur la
      façon de suivre les composants venus de l'extérieur.
\item Les équipes ont donc dû chercher pendant des semaines avant de pouvoir corriger.
\item La commission qualifie la faille de durablement présente et estime qu'on la retrouvera
      pendant une décennie.
\end{itemize}

\subsection*{10. Petite banque britannique, 2015}
\textit{Sanctions conjointes des deux autorités de contrôle britanniques.}
\begin{itemize}[itemsep=2pt]
\item Les autorités relèvent que les systèmes et les contrôles de la banque étaient insuffisants
      pour surveiller et encadrer ses contrats d'externalisation.
\item Elles notent qu'il s'agit d'un manquement déjà signalé auparavant et resté sans correction.
\item Une panne chez le prestataire qui traitait les paiements par carte a interrompu le service
      aux clients pendant plusieurs heures.
\item Le rétablissement et l'information des clients ont été jugés tardifs.
\end{itemize}

\newpage
\section*{La grille}

Une ligne par incident. Dans chaque case, écrivez le nom du domaine qui a lâché en premier, ou
\emph{en même temps}, ou \emph{non concerné}, ou \emph{je ne sais pas}.

\vspace{3mm}
\begin{center}\scriptsize
\renewcommand{\arraystretch}{2.7}
\setlength{\tabcolsep}{2pt}
% DIX COLONNES A 1,42 cm DEBORDAIENT DE 17 pt, soit 6 mm hors du bloc de texte. Le defaut
% etait la avant l'extension a dix recits. 1,36 cm rend exactement les 6 mm manquants.
\begin{tabular}{@{}p{2.05cm}|p{1.36cm}|p{1.36cm}|p{1.36cm}|p{1.36cm}|p{1.36cm}|p{1.36cm}|p{1.36cm}|p{1.36cm}|p{1.36cm}|p{1.36cm}@{}}
\toprule
\textbf{Incident} & Gouv. \newline / Incid. & Gouv. \newline / Tests & Gouv. \newline / Prest. &
Gouv. \newline / Partage & Incid. \newline / Tests & Incid. \newline / Prest. &
Incid. \newline / Partage & Tests \newline / Prest. & Tests \newline / Partage &
Prest. \newline / Partage \\
\midrule
1. Messagerie & & & & & & & & & & \\ \hline
2. Banque UK & & & & & & & & & & \\ \hline
3. Renseign. & & & & & & & & & & \\ \hline
4. Courtage & & & & & & & & & & \\ \hline
5. Banque US & & & & & & & & & & \\ \hline
6. Éditeur & & & & & & & & & & \\ \hline
7. Transfert & & & & & & & & & & \\ \hline
8. Sécurité & & & & & & & & & & \\ \hline
9. Composant & & & & & & & & & & \\ \hline
10. Cartes & & & & & & & & & & \\
\bottomrule
\end{tabular}
\end{center}

\vspace{4mm}
\textbf{Une dernière question, la plus utile.} Y a-t-il une case où vous avez vraiment hésité, ou
un enchaînement que vous avez vu sur le terrain et qui ne ressemble à aucun de ces sept récits ?

\vspace{12mm}\hrule\vspace{1mm}\hrule\vspace{1mm}\hrule

\vspace{4mm}
\begin{cle}
\textbf{Pourquoi on vous le demande.} J'ai déjà rempli cette grille de mon côté. Je ne vous
montre pas mes réponses avant que vous ayez rendu les vôtres, pour ne pas vous influencer :
l'intérêt est justement de savoir si deux personnes lisant les mêmes faits y voient le même
enchaînement. Un désaccord entre nous n'est pas un échec, c'est le résultat le plus instructif
que cet exercice puisse produire.
\end{cle}

\vspace{-1pt}
{\footnotesize\textcolor{navy}{\textbf{Merci.}} Kélian Kaddouri \textbullet{}
kkaddouri@nexialog.com}

\end{document}
